\subsection{Conventions and orientations}
The finite structures below are unlabelled structures on labelled finite sets;
all identifications are equivariant under relabelling. A statement that two
such structures are equivalent means an equivalence of their groupoids,
not merely equality of their automorphism-group orders. We distinguish
three two-element sets attached to a square: its directions of rotation,
its two diagonals, and its two partitions into opposite sides. This
convention is essential on 119/156--157.

\begin{saresult}{E.001}{Proved}{Contracted products and orientations}
For two two-element sets $P,Q$, put
$P\wedge Q=(P\times Q)/((p,q)\sim(\bar p,\bar q))$.
This is again a two-element set. Contracted product is associative and
commutative up to its canonical identifications, has the pointed torsor
$\{\pm1\}$ as unit, and $P\wedge P$ is canonically pointed by the diagonal.
For a finite set $S$ of cardinality at least two, its orientation torsor is
$\omega(S)=\operatorname{Bij}(\{1,\ldots,|S|\},S)/\mathfrak A_{|S|}$.
For smaller cardinalities use the trivial two-element torsor. There is a
canonical identification
\[
 \omega(S\amalg T)\simeq\omega(S)\wedge\omega(T)
\]
when the order of the two summands is specified. Interchanging the summands
changes this identification by $(-1)^{|S||T|}$. The permutation representation
on the basis $S$ identifies this torsor with the orientations of $\mathbb R^S$,
and with the nonzero elements of the determinant line of $\mathbb F_3^S$.
For a nondegenerate orthogonal space in characteristic different from two,
orientations mean the determinant double cover of the orthogonal frame
space; an orthonormal basis indexed by $S$ has this same permutation torsor.
\end{saresult}
\begin{proof}
The simultaneous involution on $P\times Q$ is free and has two orbits.
For three factors, either association is the quotient by the subgroup of
$(\mathbb Z/2)^3$ of even coordinate sum. This gives the canonical associator
and proves its coherence for any number of factors: all parenthesizations
have the same even-flip quotient. Swapping factors gives the symmetry.
The diagonal and off-diagonal in $P\times P$ are the two canonical classes.
Concatenating ordered listings of $S$ and $T$ gives the displayed orientation
map; changing either listing by an odd permutation changes exactly its
contracted-product class. Moving the first block past the second uses
$|S||T|$ transpositions, proving the sign. Determinants of permutation
matrices are their signs, proving the real and $\mathbb F_3$ statements.
For orthogonal frames the determinant is $\pm1$; a reflection in a vector
of nonzero norm has determinant $-1$, so both fibres occur. No order on an
arbitrary field, and no canonical orientation of its whole general linear
frame space, is asserted.
\end{proof}

\begin{saresult}{E.002}{Proved}{The small exceptional isomorphisms and their precise hypotheses}
The groupoid descriptions in 119/160 have the following realizations:
\[
\begin{aligned}
\mathfrak S_2&\simeq\mathbb F_3^*,\\
\mathfrak S_3&\simeq\operatorname{GL}_2(\mathbb F_2)
 =\operatorname{SL}_2(\mathbb F_2)
 =\operatorname{PGL}_2(\mathbb F_2)
 \simeq\operatorname{AGL}_1(\mathbb F_3),\\
\mathfrak S_4&\simeq\operatorname{AGL}_2(\mathbb F_2)
 \simeq\operatorname{PGL}_2(\mathbb F_3)
 \simeq\operatorname{Aut}^+(\text{cube}),\\
\mathfrak A_5&\simeq\operatorname{PGL}_2(\mathbb F_4),\qquad
\mathfrak S_5\simeq\mathrm{P}\Gamma\mathrm{L}_2(\mathbb F_4)
 \simeq\operatorname{PGL}_2(\mathbb F_5).
\end{aligned}
\]
The last isomorphism uses the exceptional transitive action on six points;
it is not the natural action on five points. The unqualified group of
linear maps preserving three independent lines in $\mathbb F_3^3$ has
order $48$, not $24$. Its determinant-one subgroup, or its quotient by
$\{\pm I\}$, is $\mathfrak S_4$. Thus that clause of 119/160 needs a
restriction or a quotient.
\end{saresult}
\begin{proof}
On a three-element set, adjoining zero gives a two-dimensional vector space
over $\mathbb F_2$: the sum of two different nonzero elements is the third.
Every permutation respects this rule. On a four-element set, any choice
of zero gives the same affine plane over $\mathbb F_2$, because the sum of
all four points is zero. The affine group has $4\cdot6=24$ elements and
acts faithfully. An affine line over $\mathbb F_3$ has $3\cdot2=6$
automorphisms, acting faithfully on three points. A fractional linear map
on a projective line is determined by the images of three distinct points;
$|\operatorname{PGL}_2(q)|=q(q^2-1)$ follows by dividing
$(q^2-1)(q^2-q)$ by $q-1$. These facts prove the projective assertions for
$q=2,3$.

For $q=4$, the group on five points has order $60$ and is contained in
$\mathfrak A_5$: translations have cycle type $2^2$, nonzero scalings
have type $3$, and inversion has type $2^2$; these generate the fractional
linear group by elementary matrices. Frobenius swaps the two elements
of $\mathbb F_4\setminus\mathbb F_2$ and fixes the other three projective
points, so adjoining it gives $\mathfrak S_5$. In particular its action on
$\omega(\mathbb P^1(\mathbb F_4))$ is its action on the two roots of
$X^2+X+1$. Equivariant identifications of these torsors exist; choosing one
root and one orientation fixes one, and there are two choices.

The rotational cube group permutes its four body diagonals faithfully.
It has $3!\,2^2=24$ elements, namely signed permutation matrices of
determinant one. The full signed permutation group has order $48$ and
central inversion gives its direct-product decomposition with this
rotation group. The identical matrix count over $\mathbb F_3$ proves the
correction concerning three independent lines. Finally $\operatorname{PGL}_2(5)$ has an explicit faithful action on five
objects: take the five Klein four subgroups in its even subgroup of order
sixty and act by conjugation. The attached program constructs all fractional
linear permutations, all fifteen even involutions, all commuting pairs
and hence all five Klein four subgroups. It verifies that the resulting
five-point action has 120 distinct permutations. Thus its image is all
of $\mathfrak S_5$, proving the last isomorphism without assuming an
icosahedral group identification.
\end{proof}

\subsection{Hexagrams, icosahedra, and the square twist}
\begin{saresult}{E.003}{Proved}{The outer duality, its dictionary, and a bihexagon construction}
Let $S$ have six elements. A duad is a two-element subset, a syntheme a
partition into three duads, and a total a partition of all fifteen duads
into five synthemes. Let $S^*$ be the set of totals. Then $|S^*|=6$, the
natural homomorphism $\mathfrak S(S)\longrightarrow\mathfrak S(S^*)$ is an
outer isomorphism, and $S\longrightarrow S^{**}$ is a canonical bijection.
The duality interchanges duads and synthemes, and gives a canonical
bijection between the ten unordered $3+3$ partitions on either side.
For each $s\in S$, $t\in S^*$ there is a unique
$\operatorname{Stab}(s,t)$-equivariant bijection
$S\setminus\{s\}\longrightarrow S^*\setminus\{t\}$.

Here is an explicit completion of the undefined marginal term
``bihexagones'' on 119/155. A hexagon is a Hamiltonian cycle in $K_S$,
and a bihexagon is a pair of edge-disjoint hexagons. The graph whose
vertices are hexagons and whose edges are bihexagons has six connected
components, each a Petersen graph. Its components are canonically $S^*$.
Thus the duality can indeed be constructed from bihexagons. This supplies
a precise construction; the one-word margin does not specify a unique
historical definition of that word.
\end{saresult}
\begin{proof}
There are $15$ synthemes. Two disjoint synthemes form an alternating
hexagon. Label its vertices cyclically $0,1,\ldots,5$. Among the synthemes
using no hexagon edge, one consists of the three long diagonals; the
other three each consist of one long diagonal and the two disjoint short
diagonals perpendicular to it. Only the latter three together cover the
nine remaining edges. Hence two disjoint synthemes extend to a unique
total. Each syntheme has eight disjoint partners, so counting ordered
pairs gives $15\cdot8/(5\cdot4)=6$ totals. The six totals, with every
entry explicit, appear in \texttt{logs/check\_E\_finite.txt}; the program
constructs them from all $\binom{15}{5}$ candidates, not from a presumed
list.

Two different totals share exactly one syntheme. This can be seen by
fixing the standard alternating hexagon just used and listing the two
completions containing any chosen syntheme; equivalently the six displayed
totals give all fifteen intersections. Each syntheme belongs to two
totals, so it gives a dual duad. The three synthemes containing a fixed
duad give a partition of $S^*$ into three pairs. These are exactly the
fifteen dual synthemes. For a point $s$, the five duads through $s$ give
five disjoint dual synthemes: every original syntheme contains exactly
one such duad. They therefore give a total on $S^*$. The resulting six
double-dual totals are distinct and exhaust $S^{**}$. All constructions
commute with bijections of $S$. An element fixing all totals consequently
fixes their intersections, then all duads, then all points. The action
on $S^*$ is faithful and hence an isomorphism. Applying $(01)$ to the
six explicit totals gives three transpositions; it cannot be inner.

For a triple $A\subset S$, the dual synthemes of its three duads form
$K_{3,3}$ on $S^*$. The two parts of this graph give the required dual
$3+3$ partition. The complementary triple gives the same bipartition.
It suffices to check $A=\{0,1,2\}$ in the explicit list, since the
symmetric group is transitive on triples; the program additionally checks
all ten partitions and verifies uniqueness under their order-$72$
stabilizers. The incidence relation in item 3b of 119/155 is explicitly
\[
 R(d,\{T,U\})\quad\Longleftrightarrow\quad d\in T\cap U.
\]
Its rows are the dual synthemes. It recovers the complete duality, so it
also states the missing incidence axioms without an existential choice.

The stabilizer of $(s,t)$, computed directly from the six totals, is
$\mathbb F_5\rtimes\mathbb F_5^*$ of order $20$. Its two actions on the
remaining five points are transitive and faithful; their point stabilizers
are its five conjugate cyclic subgroups of order four. Match a point in
either action with its stabilizer. This gives the asserted bijection.
A cyclic complement normalizes only itself, so there is exactly one.
Independently of the group identification, the program tests all $36$
pairs $(s,t)$ and all $5!$ proposed bijections for each pair, finding
exactly one each time. This is a complete finite proof of the last
assertion directly from the total list, with no dependence on E.004.

Finally a hexagon is exactly a pair of disjoint synthemes, its alternating
edge classes. Two edge-disjoint hexagons provide four disjoint synthemes;
the remaining three edges form the fifth. Thus both hexagons lie in the
same total. Within a total, hexagons are its ten two-element subsets of
synthemes, and disjointness of hexagons is disjointness of these subsets.
This is the graph $\operatorname{KG}(5,2)$, the Petersen graph. It is
connected: intersecting two-element subsets have a common disjoint
neighbour unless their union has four elements, in which case they were
already disjoint; identical subsets need no path. Consequently the six
totals are exactly the six components. The independent exhaustive check
finds $60$ hexagons, $90$ bihexagons, and six components of ten vertices.
The usual syntheme construction is also given in Cameron
\cite[\S2, pp.2--3]{Cameron2015}; the present bihexagon formulation is
proved above rather than attributed to that source.
\end{proof}

\begin{saresult}{E.004}{Proved}{Twelve icosahedra, six biicosahedra, and their orientations}
For $S=\{s,c_0,\ldots,c_4\}$, let $F(s;c)$ consist of the five triples
$\{s,c_i,c_{i+1}\}$ and the five complements in $\{c_0,\ldots,c_4\}$
of the diagonals $\{c_i,c_{i+2}\}$, with indices modulo five.
There are twelve such ten-face complexes on $S$, independent of which
vertex is used as $s$. Call them icosahedral structures. Complementation
inside $\binom S3$ pairs them into six biicosahedral structures.
Their automorphism groups are respectively the transitive
$\mathfrak A_5$ and $\mathfrak S_5$ in $\mathfrak S_6$.
There is a natural bijection $\operatorname{Bic}(S)\simeq S^*$, and the
explicit convention
\[
 \alpha(F(s;c))=[s,c_0,c_1,c_2,c_3,c_4]\in\omega(S)
\]
gives a natural bijection
$\operatorname{Ic}(S)\simeq S^*\times\omega(S)$.
The geometric complex is the antipodal quotient of an ordinary
icosahedron, usually called a hemiicosahedron; the adjective ``gauche''
is interpreted in this precise six-vertex sense.
\end{saresult}
\begin{proof}
The dihedral group of a pentagon has order ten and consists of even
permutations of its five vertices. Thus the $5!/10=12$ pentagons at a
fixed root give twelve different complexes and a well-defined displayed
orientation. Each edge belongs to two faces, each vertex has a pentagonal
link, each four-element subset contains two faces, and exactly one of
any pair of complementary triples is a face. These assertions follow
by substituting the indices $0,\ldots,4$ in the formula. For a fully
explicit model the ten faces are
\[
012,015,023,034,045,124,134,135,235,245.
\]
At root $1$ the link order is $(0,2,4,3,5)$. The permutation taking
$(0,1,2,3,4,5)$ to $(1,0,2,4,3,5)$ is even, and substitution in the
same formula recovers the same ten faces. The root stabilizer is
transitive on the other five vertices, so this one calculation proves
root independence and orientation independence for every vertex.
Replacing the pentagon by its complementary star gives the complementary
face set and reverses $\alpha$.

Given the pair of complementary pentagons at $s$, pair each edge of one
with the unique disjoint edge of the other; add the duad joining $s$ to
the unused vertex. The resulting five synthemes form a total. Substitution
in the displayed model proves that this total is independent of $s$;
transitivity and relabelling then prove the assertion for all structures.
Its image has six elements, so this is the required bijection to $S^*$.
This is exactly the pentagon-to-pentad construction of Howard--Millson--
Snowden--Vakil \cite[\S\S1.2--1.4, pp.2--3]{HMSV2008}; our code checks
all $72$ choices of structure and root independently.

The stabilizer of a face complex has order $720/12=60$. It preserves
$\alpha$ and hence lies in $\mathfrak A_6$. The stabilizer of its
biicosahedron is, by E.003, a point stabilizer in the outer action,
therefore is $\mathfrak S_5$; its even subgroup is $\mathfrak A_5$.
The two complexes over a biicosahedron have opposite orientations, proving
the product formula. There are abstractly two equivariant choices of
this product formula; the displayed root convention selects one.

For the geometric assertion take the twelve vectors
$(\pm1,\pm\varphi,0)$, $(0,\pm1,\pm\varphi)$,
$(\pm\varphi,0,\pm1)$, where $\varphi^2=\varphi+1$ and $\varphi>1$.
Their convex hull is the regular icosahedron. Identifying opposite
vertices gives six labels; its twenty triangular faces come in ten
opposite pairs, giving the displayed ten triples after labelling a
vertex link. Conversely the formula reconstructs them from that link.
This coordinate realization and its opposite under
$\varphi\mapsto1-\varphi$ are also explicitly recorded in
\cite[\S1.3, p.2]{HMSV2008}.
\end{proof}

\begin{saresult}{E.005}{Proved}{The three perspectives of the icosahedron}
An icosahedron with a vertex is a pentagon on the other five points;
a biicosahedron with a vertex is a complementary pair of pentagons.
A biicosahedron with a distinguished edge $\varepsilon=\{a,b\}$ is
exactly a square $Q$ on the other four vertices, together with an
identification $\varepsilon\simeq\operatorname{rot}(Q)$. Choosing one
of its two icosahedra is choosing one of the two diagonals of $Q$.
An icosahedron with a distinguished face is exactly an ordered pair of
three-element sets with a bijection between them. A biicosahedron with
a distinguished biface is the corresponding unordered prism.
The stabilizers are respectively $D_{10},\operatorname{AGL}_1(5)$,
$D_8$, $\mathfrak S_3$, and $\mathfrak S_3\times C_2$.
\end{saresult}
\begin{proof}
The vertex assertions and their inverses are the defining formula in
E.004. For the edge perspective, if $F$ is one complex, let $D$ be the
two points $x$ for which $abx$ is a face. The other complex gives the
complementary pair $D'$. These are the diagonals of a square $Q$.
For a chosen $D$, let $C_a$ be the opposite-side partition consisting
of the two cross-diagonal pairs $xy$ for which $axy$ is a face. The
analogous $C_b$ is the other opposite-side partition. A rotation of a
square is canonically a contracted-product class $[D,C]$: take a vertex
on $D$ and its successor, and let $C$ be the opposite-side pair they
determine. Choosing the other vertex of $D$ changes nothing; replacing
$D$ by $D'$ changes $C$ to the other partition. Hence
$a\mapsto[D,C_a]$, $b\mapsto[D,C_b]$ is independent of which complex
was chosen and identifies $\varepsilon$ with the rotation torsor.

Conversely, given $Q$, this identification, and a selected diagonal $D$,
recover $C_a,C_b$ by that rule. The ten faces are: $abx$ for $x\in D$;
$aD'$ and $bD'$; $a$ joined to each of the two sides in $C_a$;
$b$ joined to each of the two sides in $C_b$; and $D$ joined to each
of the two vertices in $D'$. This is a complete inverse, verified by
substituting the standard square; changing $D$ complements the face
set. There are $3\cdot2=6$ such biicosahedra for any fixed edge.
The automorphisms are the eight square symmetries, with the edge action
forced by their rotation character.

For the face perspective, let $f$ be a face and $f'=S\setminus f$.
Match $x\in f$ with the unique $x'\in f'$ such that
$(f\setminus\{x\})\cup\{x'\}$ is a face. For the displayed model and
$f=012$ the matching is $0\mapsto4$, $1\mapsto3$, $2\mapsto5$.
It is a bijection. Given any such matching, the inverse face set consists
of $f$ itself, the three triples $(f\setminus\{x\})\cup\{x'\}$, and
the six triples $\{x,x',y'\}$ with $x\ne y$ in $f$. These ten triples
are exactly E.004's model after relabelling. Swapping the two matched
triads complements the complex, so forgetting the choice gives precisely
the unordered prism and its $\mathfrak S_3\times C_2$ symmetry.
The face stabilizer is $\mathfrak S_3$. The program exhausts all $90$
biicosahedron-edge pairs and all $120$ icosahedron-face pairs and verifies
these reconstruction maps.

The orientations from the three perspectives agree with E.004 after
one explicit sign convention. In the face perspective it is minus the
orientation of $(x,y,z,x',y',z')$. In the edge perspective choose
$u\in D$, follow the rotation attached to $a$ to write
$(u,v,u',v')$, and take minus the orientation of
$(a,b,u,v,u',v')$. Changing a permissible choice is even; comparison
in the displayed ten-face model gives the sign. Thus these are complete
versions of the three proposed definitions of $\alpha(F)$.
\end{proof}

\begin{saresult}{E.006}{Proved}{Square torsors and the corrected twist}
For a square $Q$, write $D(Q)$ for its two diagonals, $C(Q)$ for its two
opposite-side partitions, and $W(Q)$ for its two rotation directions.
Then
\[
 W(Q)=D(Q)\wedge C(Q),\qquad\omega(Q)=C(Q)
\]
where the second $\omega$ is the sign orientation of the underlying
four-element set. For a square on four of six points, let
$\varepsilon$ be its complementary pair. Under E.003 its stabilizer
preserves a unique square $Q^*$ and complementary pair $\varepsilon^*$
on $S^*$. Their torsors satisfy
\[
 D(Q^*)=C(Q),\quad C(Q^*)=D(Q),\quad
 W(Q^*)=W(Q),\quad
 \varepsilon\wedge\varepsilon^*\wedge W(Q)=1.
\]
Let $Q^\vee$ be the square of sides of $Q$. Define
\[
 \varepsilon\wedge Q^\vee
 = (\varepsilon\times Q^\vee)/((e,c)\sim(\bar e,\bar c)),
\]
where $\bar c$ is the opposite side. There are exactly two equivariant
square isomorphisms
$Q^*\simeq\varepsilon\wedge Q^\vee$, differing by the antipodal map.
There is no equivariant isomorphism to the analogous twist of the
original vertex square $Q$. Thus 119/157's displayed dual is essential.
The formula on 156 is true if its notation $\varepsilon\wedge Q$
is defined to include this side dual. Reading its $\omega(Q)$ as the
sign orientation of four vertices makes its first torsor equation false.
\end{saresult}
\begin{proof}
The dihedral group acts on $D,C,W$ by three characters, with
$W=DC$, as the successor construction in E.005 shows. A quarter-turn
exchanges both diagonals and opposite-side partitions; a reflection
through opposite vertices fixes $D$ and exchanges $C$. These generate
$D_8$, and their signs on four vertices are respectively $-1,-1$.
Thus the sign character is $C$, not $W$.

Here is a fully explicit finite proof of the dual assertions. Take
$Q=(0,1,2,3)$ cyclically, $\varepsilon=\{4,5\}$, and number totals
lexicographically as in the program's output. The order-$16$ stabilizer
is generated by $(0123)$, $(13)$, and $(45)$. Its two orbits on totals
are $\varepsilon^*=\{2,4\}$ and $Q^*=\{0,1,3,5\}$. The unique invariant
pairing on the latter is its square structure. For these three generators,
or equivalently for all sixteen elements as checked in the program,
the characters are exactly the four displayed relations. Checking generators
suffices because each is a homomorphism to $\{\pm1\}$.

Represent a class of the proposed twist by $(4,c_i)$, with sides
$c_0=01,c_1=12,c_2=23,c_3=03$ in cyclic order. If $g$ exchanges $4,5$,
its side action is followed by the antipodal shift $i\mapsto i+2$.
The assignments
\[
 (c_0,c_1,c_2,c_3)\longmapsto(0,5,1,3)
 \quad\hbox{or}\quad(1,3,0,5)
\]
intertwine the three generators with their explicit action on totals.
They are therefore equivariant square isomorphisms. Any such map is
determined by the image of $c_0$; its stabilizer has precisely two fixed
vertices in $Q^*$, proving that these are the only two. By contrast a
central twist of the vertex square leaves its diagonal and opposite-side
characters unchanged, while duality exchanges them. Since those characters
are different, no vertex-square isomorphism exists. The program also
checks all $4!$ bijections in both cases, giving exactly two and zero.

Finally dualizing twice returns $Q$. Indeed $Q^{\vee\vee}\simeq Q$,
$\varepsilon\wedge\varepsilon^*=W(Q)$, and there is a canonical
isomorphism $W(Q)\wedge Q\simeq Q$: send $[w,q]$ to the next vertex
from $q$ in direction $w$. Reversing $w$ and replacing $q$ by its
antipode does not change that next vertex. This also proves the reciprocal
twist formula, with the same convention concerning side duals.
\end{proof}

\subsection{The affine plane and triality}
\begin{saresult}{E.007}{Proved}{The incidence axioms on 119/159, with two corrections}
A set of nine points carrying twelve three-element blocks, any two
\emph{distinct} blocks meeting in at most one point, is the affine plane
$\mathbb F_3^2$, uniquely up to isomorphism. The quantifier must exclude
equal blocks. Dually, partition a twelve-element set as
$E=\coprod_{i=1}^4E_i$, with $|E_i|=3$. Conditions a) and c) on 159,
namely four-element transversals and unique blocks through any pair
from different classes, imply that there are nine blocks, that every
point belongs to three, and that any two different blocks meet in
exactly one point. They characterize the dual affine plane. Adjoining
$\infty$ and the four blocks $E_i\cup\{\infty\}$ gives
$\mathbb P^2(\mathbb F_3)$.
Condition c) alone does not imply nine blocks, even if every block has
four elements.
\end{saresult}
\begin{proof}
The twelve triples contain $12\binom32=36=\binom92$ pairs and no pair
repeats, so every pair occurs once. Each point belongs to four lines.
Given a line $L$ and a point $p\notin L$, the three joins of $p$ to the
three points of $L$ are different. Its fourth line is the unique line
through $p$ disjoint from $L$. The disjoint lines partition the six
points outside $L$ into two triples. Thus disjointness gives four
parallel classes, each of three lines.

Use two classes as rows and columns of a $3\times3$ array. Label their
coordinates by $\mathbb F_3$. A third class through $(0,0)$ can, after
interchanging labels $1,2$, be made to contain $(1,1),(2,2)$. Its other
two lines are the two derangements of three letters, hence are
$y=x+1$ and $y=x-1$. The fourth class must use the remaining pairs,
and is $y=-x+c$. This proves the claimed uniqueness and gives explicit
coordinates.

For the dual axioms there are $\binom42\cdot3^2=54$ cross-class pairs,
and every transversal contains six of them, so c) gives $9$ blocks.
At a point there are nine cross-class partners, three in each incident
block, so there are three incident blocks. Two blocks meet in at most
one point by c). A fixed block meets $4(3-1)=8$ distinct other blocks,
so it meets every other block in exactly one point. Reversing incidence
now gives the first design. Completing the parallel classes by one
point is the standard projective completion of its dual: in homogeneous
coordinates it is the incidence of one- and two-dimensional subspaces
of $\mathbb F_3^3$. It has thirteen points and thirteen lines.

For the final assertion order the four classes. For every $i<j$ and
every $b\in E_j$, take the block $E_i\cup\{b\}$. There are eighteen
different four-element blocks. A cross-class pair $a\in E_i,b\in E_j$,
$i<j$, occurs in exactly $E_i\cup\{b\}$. Thus c) holds, while the
claimed cardinality fails. The program verifies all $54$ cross-class
pairs of this completely specified example. It fails a), which is
precisely the missing hypothesis in the marginal inference.
\end{proof}

\begin{saresult}{E.008}{Proved under stated hypotheses}{The three symmetric groups and the nonsplit affine extension}
Let $V=\mathbb F_3^2$, let $E$ be its twelve lines, and let $I$ be its
four directions. Then $\operatorname{Aut}(V,E)=\operatorname{AGL}_2(3)$
has order $432$, and
\[
1\longrightarrow V\rtimes\{\pm I\}
 \longrightarrow\operatorname{AGL}_2(3)
 \longrightarrow\mathfrak S_4\longrightarrow1
\]
is nonsplit. The kernel is the equal-sign subgroup of
$\mathfrak S_3\times\mathfrak S_3$ on any two chosen parallel classes.

For each pairing $\alpha$ of the four directions, write
$E=A_\alpha\amalg B_\alpha$ for the resulting two six-line sets.
There is exactly one graph of an isomorphism
$\mathfrak S(A_\alpha)\longrightarrow\mathfrak S(B_\alpha)$ that
contains the translations and is normalized by the block-preserving
part of $H_\alpha=\operatorname{Stab}_{\operatorname{AGL}_2(3)}(\alpha)$.
It is the graph of an outer isomorphism; denote it by $G_\alpha$.
The whole affine group permutes these three groups. Moreover
\[
H_\alpha\simeq\mathrm{A}\Gamma\mathrm{L}_1(9),\quad
|H_\alpha|=144,\quad |H_\alpha\cap G_\alpha|=72,
\]
and conjugation embeds $H_\alpha$ with index ten in
$\operatorname{Aut}(G_\alpha)$, of order $1440$. The translations are a
common Sylow $3$-subgroup, and
\[
N_{G_\alpha}(V)/V\simeq D_8,\qquad
N_{\operatorname{Aut}(G_\alpha)}(V)/V\simeq\operatorname{SD}_{16}.
\]
Thus ``order ten'' on 156 must be index ten; and the square-group
normalizer assertion is true with $G_\alpha$, not with
$\operatorname{Aut}(G_\alpha)$. Without the normalization condition
there are nine translation-compatible outer graphs for each pairing,
so an unqualified uniqueness assertion is false.
\end{saresult}
\begin{proof}
An incidence automorphism is determined by the images of an origin and
an ordered affine basis. There are $9\cdot8\cdot6=432$ choices, each
realized by one affine map. Fixing the four directions forces its linear
part to be scalar; the kernel is $V\rtimes\{\pm I\}$ of order eighteen.
On two parallel classes it acts as $(x,y)\mapsto(ax+b,ay+c)$,
$a=\pm1$, giving exactly the equal-sign subgroup. The quotient is
$\operatorname{PGL}_2(3)=\mathfrak S_4$ by E.002.

If a linear lift of a double transposition on directions had square
$I$, it would be diagonalizable over $\mathbb F_3$, and would fix a
direction. Therefore such a lift has square $-I$. Any affine lift
$(M,v)$ has order four, because $I+M+M^2+M^3=0$ and its square has
linear part $-I$. A section would lift a double transposition to an
involution, which is impossible. This also explains the exact check:
all $54$ lifts of the three double transpositions have order four.

We give an exhaustive finite construction of the three graphs, so their
existence and uniqueness require no classification of unknown subgroups
of $\mathfrak S_{12}$. Number directions
$(1,0),(0,1),(1,1),(1,2)$, and number the three lines in each direction
lexicographically by their three point coordinates. Number totals as in
E.003, and write $\Phi$ for its outer permutation isomorphism. All
isomorphisms between two six-element permutation groups have one of the
forms
\[
 x\longmapsto fxf^{-1},\qquad
 x\longmapsto\Phi^{-1}(fxf^{-1}),\qquad f\in\mathfrak S_6.
\]
Here is a proof that this list is exhaustive. An automorphism sends a
transposition to an involution in a conjugacy class of size fifteen;
the only such classes are transpositions and triple transpositions.
If it preserves transpositions, it permutes the maximal five-element
families of pairwise noncommuting transpositions. Such a family is the
star at one point: pairwise intersecting two-element subsets are either
a star or a three-edge triangle. It therefore permutes the six points,
and is inner. Composing with $\Phi$ reduces the other case to this one.
Distinct $f$ give distinct maps because the centre is trivial.

For each of the three pairings, the program enumerates these $1440$
possibilities. It tests equality on all nine translations, then tests
normalization for every element of the order-$72$ block-preserving
stabilizer and the two generators $(01),(012345)$ of $\mathfrak S_6$.
Testing generators suffices because both sides of the intertwining
identity are homomorphisms in $x$. Precisely nine outer maps survive
the first test, and one survives the second; no inner map survives even
the first. The surviving permutations $f$, in the stated deterministic
numbering, are
\[
\begin{array}{c|c}
\alpha&f\text{ as the list }(f(0),\ldots,f(5))\\\hline
01\mid23&(3,0,4,5,2,1)\\
02\mid13&(5,1,2,3,4,0)\\
03\mid12&(3,4,0,5,1,2).
\end{array}
\]
These formulas specify all $720$ elements of each $G_\alpha$ and
are a finite certificate independent of a group-name recognition routine.
The code explicitly checks that all $432\cdot3$ conjugates are one of
these three groups, that each $H_\alpha$ normalizes its graph, and that
its intersection has order $72$. It checks that only its identity
centralizes all graph elements. Thus the asserted injection and index
follow from $|\operatorname{Aut}(\mathfrak S_6)|=2\cdot720$, proved
by the preceding classification.

The linear part of $H_\alpha$ is the inverse image of a square stabilizer
$D_8\subset\operatorname{PGL}_2(3)$, hence has order sixteen. Identify
$\mathbb F_9=\mathbb F_3[i]$, $i^2=-1$. Multiplication by $1+i$ has
order eight; Frobenius has order two and conjugates it to its cube.
Their projective images preserve one pairing of the four directions,
and their order-sixteen group is its full inverse image. This proves
$H_\alpha=\mathbb F_9\rtimes(\mathbb F_9^*\rtimes\operatorname{Gal}
(\mathbb F_9/\mathbb F_3))$, and its quotient by translations is
$\operatorname{SD}_{16}=\langle r,s:r^8=s^2=1,srs=r^3\rangle$.

In a six-point symmetric group the normalizer of a Sylow group
$C_3^2$, acting on two triples, is
$(\mathfrak S_3\times\mathfrak S_3)\rtimes C_2$ of order $72$.
The quotient by $C_3^2$ is $C_2\wr C_2=D_8$. There are ten Sylow
$3$-subgroups, so the normalizer in its full automorphism group has
order $1440/10=144$. The already embedded $H_\alpha$ normalizes the
translations and has this order, so it is exactly that normalizer.
This proves both normalizer readings, including the distinction omitted
in R12.
\end{proof}

\begin{saresult}{E.009}{Refuted}{The literature-absence claim L7}
The mathematical presentation on 119/156 has an explicit counterpart
in the literature. Cameron's 2015 account of the dual Witt designs
contains the same affine plane, the same action on its twelve lines,
and the same three outer $\mathfrak S_6$ graph subgroups. Consequently
L7, interpreted as absence of this construction from the literature on
the outer automorphism, is false. This does not assert that Cameron
studied Grothendieck's manuscript or that the terminology agrees.
\end{saresult}
\begin{proof}
The source has been opened and its construction read, not inferred
from a search snippet. Cameron \cite[\S2, pp.2--3]{Cameron2015}
constructs the outer action by totals. Section 3, p.4, constructs
$S(5,6,12)$ on two dual six-element sets. Section 4, p.5, proves that
the stabilizer of a hexad acts as $\mathfrak S_6$ on it and by the outer
action on its complement. Section 5, pp.5--8, constructs the dual pair
of twelve-element Witt designs. Finally \S6, pp.8--9, identifies the
stabilizer of a three-element set in one design with
$\operatorname{AGL}_2(3)$ on its remaining nine points, and identifies
the twelve points of the dual design with the lines of that affine
plane. It identifies parallel classes with the four triangles of the
corresponding $3+3+3+3$ partition.

Here is the exact derivation of our three groups from those statements.
The three pairs in the distinguished triple correspond, by Cameron's
duality, to three complementary hexad pairs on the dual twelve-set.
Their common refinement is the four parallel classes, so these are
exactly its three pairings $\alpha$. The stabilizer of either hexad
in a complementary pair is the same subgroup, a graph of the outer
$\mathfrak S_6$ isomorphism. Every translation preserves each parallel
class, hence each hexad, and lies in this graph. The pairing stabilizer
normalizes it, including when it swaps the two complementary hexads.
The uniqueness proved in E.008 therefore identifies Cameron's three
subgroups with $G_\alpha$, not merely with three abstractly isomorphic
groups. This is a precise published-location match to the proposed
triality. DeVos \cite[pp.6--8]{DeVosBlocks} independently describes the
$\operatorname{AGL}_2(3)$ three-set stabilizer and its permutation of
the three pairings of four slopes; this second source is corroboration,
not needed to fill a gap in the first derivation.
\end{proof}

\subsection{Word types, chord diagrams, and maps}
\begin{saresult}{E.010}{Proved under stated hypotheses}{The exact meaning of word type in 152/003}
Consider an oriented disc $D$, a finite union $K$ of disjoint properly
embedded arcs with distinct endpoints on $\partial D$, a nonempty
$n$-element set $S\subset\partial D\setminus K$, and a root $s_0\in S$.
Every component of $\partial D\setminus\partial K$ must contain at least
one mark, and every such component outside the region containing $s_0$
must contain exactly one. These are the data and conditions on 003.
Their isotopy groupoid is equivalent to the following explicit groupoid
of word types: a rooted cyclically ordered $n$-set of marks, together
with a partial pairing of its $n$ gaps, noncrossing, such that after
cutting at the root no pair encloses an unpaired gap. Morphisms preserve
the root, cyclic order, and pairing.

Equivalently a type can be written as a linear word of $n$ positions
with a specified cancellation matching: paired letters are the two
occurrences of a private symbol with opposite exponents; every unpaired
letter has its own private symbol; the pairing is noncrossing and encloses
no unpaired letter. Renaming symbols and reversing the chosen name of
a paired symbol changes no type. Unspecified equality patterns of ordinary
free-group words are not the object classified by the page.

Under the reading ``isotypique'' as a set of isomorphism types, take
isomorphism classes of these same finite data. Under the reading
``isotopique'', retain their combinatorial isomorphisms, which are exactly
the isotopy classes of homeomorphisms. If orientation is forgotten,
allow order reversal; its possible reflection automorphisms must be
retained for a groupoid equivalence. Neither reading identifies the raw
groupoid of all homeomorphisms with a finite combinatorial groupoid.
\end{saresult}
\begin{proof}
A gap between consecutive marks contains at most one arc endpoint:
two endpoints there would cut off an interval containing no mark.
Thus endpoints select distinct gaps, and the arcs pair the selected gaps.
Two arcs cannot interlace their endpoints, by separation in a disc.
An unpaired gap joins two consecutive marks on one component of the
boundary minus endpoints. Such a component is permitted only in the root
region. This is equivalent to the requirement that no paired interval,
in the linear boundary order cut at $s_0$, encloses an unpaired gap.
This proves the necessity of all the stated combinatorial conditions.

Conversely put the marks in cyclic order on a round circle and put one
endpoint in each selected gap. Join paired endpoints by straight chords.
Noncrossing makes their interiors disjoint. Every complementary boundary
interval contains a mark; any interval containing more than one mark
contains an unpaired gap and hence belongs to the root region. Thus this
constructs exactly an object of the manuscript's category.

A combinatorial isomorphism extends first to an orientation-preserving
boundary homeomorphism carrying the endpoints, then over the arcs.
Cutting along the arcs gives discs, and the boundary maps on each extend
over the discs by radial extension after choosing disc coordinates.
The extensions glue, proving fullness and essential surjectivity.
Two extensions inducing the same finite data differ by a map fixing all
marked points and all arcs setwise. Isotope its boundary restrictions on
the finitely many intervals to the identity, then its restrictions to
the arcs to the identity. Cutting now leaves homeomorphisms of discs
fixed on their boundaries. Alexander's radial isotopy makes each one
the identity, and the isotopies glue. Thus the finite isomorphism records
exactly the isotopy class, proving faithfulness. The unoriented argument
first composes an orientation-reversing map with a fixed reflection.
In the oriented rooted case a cyclic-order automorphism fixing the root
fixes every mark and gap, so all combinatorial automorphism groups are
trivial. Passing to isomorphism classes therefore gives the alternative
``isotypique'' reading as a discrete category. Without orientation,
symmetric diagrams can retain a reflection, so a mere set of types loses
that automorphism.

The private-symbol encoding is equivalent to the partial matching by
construction, with the root supplying the linear order. Its cancellation
criterion is proved in E.011. Notice one literal boundary case in the
page's displayed ``reduced order'': if $K$ has $m>0$ arcs, its formula
$\sum(\#(S\cap\alpha)-1)$ equals $n-2m$; if $K=\varnothing$, it equals
$n-1$. Thus that formula cannot uniformly mean the usual reduced word
length, which would be $n$ in the all-unpaired case. The combinatorial
equivalence just proved does not silently replace this exceptional value.
Finally the raw homeomorphism groupoid has, even for a disc with one
mark and no arcs, infinitely many automorphisms supported in its interior.
The finite-data groupoid has none. This proves the stated limitation on
any literal non-isotopic interpretation.
\end{proof}

\begin{saresult}{E.011}{Proved}{The repaired deletion criterion, including fixed points}
An involution on a linearly ordered finite set is reducible to the identity
by repeatedly deleting \emph{both} members of an adjacent two-cycle if
and only if its two-cycles do not interlace and no two-cycle encloses a
fixed point. For a fixed-point-free involution this is exactly
noncrossing. Deleting only one member and making the survivor fixed,
as a literal completion of 152/004's one-edge contraction, fails already
for $(1\ 4)(2\ 3)$.
\end{saresult}
\begin{proof}
Deleting an adjacent pair preserves noninterlacing and the absence of an
enclosed fixed point, and inserting such a pair also preserves these
conditions. Thus any successful sequence implies them. Conversely,
choose a two-cycle $(a,b)$ with the fewest positions strictly between its
members. No intervening position is fixed. The partner of any intervening
position must also intervene, since otherwise its cycle interlaces
$(a,b)$. That would give a pair with fewer intervening positions,
contradicting minimality. Thus $(a,b)$ is adjacent. Delete both positions
and apply induction. If no two-cycle remains the involution is the
identity, completing the proof for arbitrary fixed points as well as
pairings.

For the nested pairing the only adjacent pair is $(2,3)$. Contracting one
of its positions leaves the other as a fixed point between positions
$1,4$, which are no longer eligible for adjacent contraction. The process
stops at a nonidentity involution. Deleting both positions instead leaves
one adjacent pair and succeeds. The program checks this counterexample
and tests the general repaired criterion on every involution through ten
positions. It separately tests every fixed-point-free pairing through
fourteen positions, including all $135135$ pairings at that size.
\end{proof}

\begin{saresult}{E.012}{Proved}{The complete map dictionary and the genus-zero criterion}
A finite oriented combinatorial map without isolated vertices is a finite
dart set $D$, a fixed-point-free involution $\alpha$ pairing its darts
into edges, and a permutation $\sigma$ whose cycles are the cyclically
ordered vertices. Its boundary permutation is
$\tau=\sigma\alpha$; its faces are the cycles of $\tau$.
Connectedness is transitivity of $\langle\sigma,\alpha\rangle$.
Conversely these permutations reconstruct the oriented ribbon
neighbourhood and, after capping its boundary components, the closed
oriented surface, uniquely up to oriented homeomorphism. With one face,
$|D|=2m$ and $c(\pi)$ denoting number of cycles,
\[
 g=\frac{m+1-c(\tau\alpha)}2.
\]
The following are equivalent: the graph is a tree; $g=0$; its boundary
pairing is noncrossing. This is exactly the corrected content of
152/010--013. The point tree is an additional zero-dart object.
\end{saresult}
\begin{proof}
Make an oriented disc for every cycle of $\sigma$, mark its dart intervals
in that cyclic order, and attach an untwisted ribbon between each pair
of intervals prescribed by $\alpha$. This constructs an oriented compact
surface with boundary retracting to the graph. Following one side of a
ribbon crosses the edge by $\alpha$ and moves to the next interval by
$\sigma$, so boundary components are precisely cycles of $\sigma\alpha$.
Capping them with discs gives a cellular embedding. Conversely cutting
a cellular embedding into vertex discs and edge ribbons recovers the
permutations. Maps commuting with both permutations extend over these
pieces; maps inducing the same permutation data are isotopic by the
same disc argument used in E.010. This proves the groupoid dictionary,
including automorphisms, rather than only a count of objects.

Euler's formula gives
$2-2g=c(\sigma)-m+c(\tau)$; in the one-face case it is the displayed
formula, since $\sigma=\tau\alpha$. A connected finite graph is a tree
if and only if $c(\sigma)=m+1$: a spanning tree has one fewer edge than
vertices, and every remaining edge creates a cycle. Thus tree and genus
zero are equivalent here.

A leaf is a one-cycle of $\sigma$. Its paired boundary sides are adjacent,
because $\sigma d=d$ implies $\tau(\alpha d)=d$. Removing the leaf and
its edge deletes those two adjacent boundary sides. Induction therefore
shows that a tree has a noncrossing pairing. Conversely a noncrossing
pairing has an adjacent pair by E.011; deleting it and reinserting the
corresponding leaf proves by induction that its graph is a tree. This
also gives a direct proof without presupposing Euler's classification of
surfaces. The program computes $\sigma=\tau\alpha$ for every pairing
through $m=7$ and compares its cycle count with noncrossing and deletion.
For $m=7$ it obtains genera $0,1,2,3$ with counts
$429,12012,66066,56628$, respectively, summing to $135135$.

The manuscript's final dictionary $\rho_\Gamma=\rho_C\sigma_C$ is
our $\sigma=\tau\alpha$. Its earlier line identifying the two rotations
without an edge flip cannot be used simultaneously. The two-vertex,
one-edge tree already distinguishes them: its vertex permutation is
identity, while its boundary permutation is the transposition. If the
opposite contour orientation is used all cyclic orders reverse and the
usual conjugate inverse convention results; it gives the same map and
genus. This treats the inverse-sign uncertainty without choosing an
inconsistent mixture of conventions.
\end{proof}

\begin{saresult}{E.013}{Known}{Relationship with the cartes of the Esquisse}
The dictionary of E.012 is the ordinary oriented map dictionary. Subdividing
each edge by a white midpoint makes the original vertices black and
turns the map into a clean dessin: every white vertex has degree two.
On the same dart set its permutations are
\[
 \sigma_0=\sigma,\qquad\sigma_1=\alpha,
 \qquad\sigma_\infty=(\sigma\alpha)^{-1},
 \qquad\sigma_0\sigma_1\sigma_\infty=1.
\]
Plane trees are precisely its connected genus-zero one-face instances.
The passage in the Esquisse to algebraic curves and Galois action is
additional arithmetic structure on this same combinatorial category,
not a different boundary-pairing criterion.
\end{saresult}
\begin{proof}
The original darts are exactly the half-edges obtained by subdividing
an edge. Rotation around a black vertex is $\sigma$; rotation around
a white midpoint exchanges the two half-edges, hence is $\alpha$.
The face calculation in E.012 proves the product relation and the genus
assertion. Conversely, from a bipartite map in which every white vertex
has degree two, suppress those vertices to recover the original map;
the constructions are inverse on objects and morphisms.

The primary Esquisse text has been opened. In \S3, printed pp.14--15
(PDF pp.10--11 in the opened edition), Grothendieck defines the isotopy
category of maps by flags and the cartographic generators, then uses
$\rho_s\rho_f=\sigma$, $\sigma^2=1$, to identify it with coverings
ramified over three points, with index one or two over the middle point.
Our fixed-point-free $\alpha$ is the no-semi-edge, index-two case of
that exact presentation. Printed pp.15--17 then add complex structures,
algebraic realization, and Galois action. This establishes the claimed
relationship by the displayed translation of generators, rather than by
terminological similarity \cite[\S3, pp.14--17]{GrothendieckEsquisse}.
\end{proof}

\begin{saresult}{E.014}{Proved}{Geometric realizations, hypercubes, and cross-ratios}
All the configurations above have explicit geometric realizations over
their indicated fields. An $n$-cube with labelled coordinate directions
is a product $\prod_{i\in I}P_i$ of two-element coordinate torsors; its
dual cross-polytope has vertices $\coprod_iP_i$. Its orientation torsor
is $\omega(I)\wedge\bigwedge_iP_i$, and its antipodal map has sign
$(-1)^n$. For $n=3$, a cube is equivalently its four body diagonals and
its orientation torsor, with automorphism group $\mathfrak S_4\times C_2$.
The affine and projective planes in E.007 are realized over $\mathbb F_3$;
they cannot be realized faithfully by nine distinct real points and
straight lines with exactly the affine-plane incidences.

The anharmonic ratios in 119/160 are the usual six transforms
\[
\lambda,\quad1-\lambda,\quad\lambda^{-1},\quad
(1-\lambda)^{-1},\quad\frac{\lambda-1}{\lambda},\quad
\frac{\lambda}{\lambda-1}.
\]
On $\mathbb P^1(\mathbb F_3)$ they all equal $-1$. On
$\mathbb P^1(\mathbb F_4)$ an ordering of its five points gives an
orientation precisely by choosing one of the two values in
$\mathbb F_4\setminus\mathbb F_2$, as stated on the page.
\end{saresult}
\begin{proof}
Choose a sign coordinate in each $P_i$. The cube becomes
$\{\pm1\}^n\subset\mathbb R^n$, and its dual is
$\{\pm e_i\}\subset\mathbb R^n$. Their incidences are given by choosing
coordinates or coordinate signs, so the construction is independent of
these choices. A signed coordinate permutation has determinant the sign
of the permutation times the product of coordinate signs. This proves
the torsor formula and the antipodal sign. In dimension three, the four
body diagonals are permuted faithfully by the rotation group, as proved
in E.002; central inversion reverses orientation and fixes all four.
Conversely, the four vertices of a regular tetrahedron centred at the
origin, together with their negatives, are the eight vertices of a cube.
This reconstructs the cube from the four labelled diagonals, with the
two choices recorded by the orientation torsor. Regular squares and
pentagons are obtained by points on the unit circle, and the icosahedral
coordinate realization was given in E.004. Taking its antipodal quotient
realizes the six-vertex complex on $\mathbb R\mathbb P^2$.

For the finite planes use points $(x,y)\in\mathbb F_3^2$ and lines
$ax+by=c$, modulo a common nonzero scalar on $(a,b,c)$ with $(a,b)\ne0$.
Projective completion uses nonzero triples modulo scalars and equations
$ax+by+cz=0$. These formulas realize every stated incidence and every
affine or projective transformation. For a general finite incidence
configuration in this chapter, its flag complex also supplies an entirely
explicit real geometric realization: place one basis vector for each
incident object in a real coordinate space and take the simplices spanned
by mutually incident flags. The incidence is recoverable as the face
relation. This last realization makes no false assertion that all finite
field lines become real straight lines.

Indeed every finite noncollinear set of real points has a line containing
exactly two of them. To prove this, minimize the positive distance from
a point $p$ to a line $L$ determined by two points of the set. If $L$
contains at least three points, two of them, $a,b$, lie on the same
closed half-line from the perpendicular foot, with $b$ nearer the foot.
In the plane spanned by $p,L$, the distance of $b$ from the line $pa$ is
strictly less than the chosen distance, by the area formula for the
triangle $pab$: its ratio is $|a-b|/|p-a|<1$. This is a contradiction.
An affine plane of order three has a third point on every pair's line;
it therefore has no faithful such real straight-line realization. The
all-collinear case cannot realize its distinct lines either.

Finally normalize three distinct projective points to $\infty,0,1$.
The fourth coordinate is $\lambda$; permuting the ordered quadruple
gives exactly the six displayed expressions by direct substitution in a
fractional linear map. For $\mathbb F_3$, the sole available value is
$2=-1$, fixed by all six transforms. For $\mathbb F_4$, the two values
are roots of $X^2+X+1$. The transforms $1-\lambda$ and $1/\lambda$
exchange them, and the alternating subgroup of four-letter permutations
fixes them. In an ordered five-tuple the fifth coordinate is the other
root; exchanging the fourth and fifth entries also exchanges the values.
These transpositions generate $\mathfrak S_5$, so the value depends
exactly on the parity class of the ordering. This is an explicit map
$\omega(\mathbb P^1(\mathbb F_4))\simeq\mathbb F_4\setminus\mathbb F_2$,
with the chosen normalization fixing the cross-ratio convention.
\end{proof}
